%% ma: L10-B14-0054
%% lop: 10
%% chuong: 5
%% bai: 14
%% dang: 10.14.9
%% loai: TL
%% muc: VD
%% dap_an: "-"
%% nguon: "chuyen-de-cac-so-dac-trung-cua-mau-so-lieu-khong-ghep-nhom-toan-10-chuong-trinh-moi.pdf (GV. Phan Nhật Linh) (P05), Bài 03 Dạng 2 Bài tập tự luận 14"
%% kien_thuc: "Lớp 10 Bài 13 (số trung bình), Bài 14 (phương sai, độ lệch chuẩn)"
%% trang_thai: chinh-thuc
%% ghi_chu: "Lời giải gốc tính sai điểm trung bình của Bình và Cường (ghi đều 87,67; đúng 89,11 và 88,33) và dùng phương sai chia n - 1; kho tính lại theo SGK (chia n), kết luận cử Bình không đổi. Tự duyệt theo yêu cầu giáo viên 22/10/2026."
\begin{ex}
Ba bạn An, Bình, Cường có bảng điểm như sau:
\begin{center}
\begin{tabular}{|c|c|c|c|c|c|c|c|c|c|}\hline
 & Toán & Lí & Hóa & Sinh & Sử & Địa & Ngoại ngữ & GDCD & Văn\\\hline
An & $95$ & $85$ & $88$ & $75$ & $92$ & $85$ & $93$ & $98$ & $78$\\\hline
Bình & $92$ & $89$ & $85$ & $85$ & $88$ & $94$ & $94$ & $90$ & $85$\\\hline
Cường & $93$ & $82$ & $95$ & $88$ & $83$ & $91$ & $91$ & $92$ & $80$\\\hline
\end{tabular}
\end{center}
a) Hãy tính điểm trung bình của mỗi bạn (vai trò của các môn học như nhau).

b) Lớp cần cử $1$ trong $3$ bạn đi thi kiến thức tổng hợp thì em đề cử bạn nào? Vì sao?
\loigiai{a) $\overline x_{An}=\dfrac{789}9\approx87{,}67$; $\overline x_{Bình}=\dfrac{802}9\approx89{,}11$; $\overline x_{Cường}=\dfrac{795}9\approx88{,}33$.

b) Phương sai: $s^2_{An}\approx52{,}89$; $s^2_{Bình}\approx12{,}10$; $s^2_{Cường}\approx25{,}78$.

Bình có điểm trung bình cao nhất và phương sai nhỏ nhất (học đều các môn nhất) nên đề cử Bình.}
\end{ex}
